\section{Introduction}\label{sec:intro}

Training a neural network on a digital processor moves every weight from memory to arithmetic and back at every update.
In 1960, Widrow and Hoff took a different route: they stored each weight of their adaptive ADALINE neuron in a
\emph{memistor}, an electrochemical cell whose conductance is set by the time integral of a current through a third
terminal~\cite{widrow1960adaline,widrow1960switching,widrow1990}. The weight was read without being disturbed, and it was
adapted in place. The device was macroscopic, the cells were wet, and no multilayer learning rule existed yet, so the line ended
with the Madaline hardware of the early 1960s~\cite{widrow1990}. That machine was a working analog learning system.

Sixty years later, three-terminal electrochemical devices have returned at the nanoscale. Lithium-ion
ECRAM~\cite{tang2018,fuller2017}, CMOS-compatible metal-oxide ECRAM~\cite{kim2019}, protonic programmable
resistors~\cite{onen2022,xu2026} and open-loop programmable arrays~\cite{chen2023} all decouple the read path from the write path,
just as the memistor did. They report hundreds to thousands of conductance levels, nanosecond pulses and femtojoule-scale write
energies. In parallel, algorithms for training on imperfect analog weights have matured: the stochastic pulse-coincidence update
of the resistive processing unit (RPU)~\cite{gokmen2016}, Tiki-Taka~\cite{gokmen2020} and its successors
TTv2~\cite{gokmen2021}, chopped TTv2 and AGAD~\cite{rasch2024}, and residual learning with convergence guarantees~\cite{wu2025}.

This paper asks whether memistor-class three-terminal electrochemical weights, from Widrow's cell to today's ECRAM, can support
analog in-memory \emph{training} at near-digital accuracy. If they can, what device specifications does that need, and what
energy advantage remains once the periphery is counted? We could not build devices. Instead we provide what groups that can
build them need to test the idea quickly:

\begin{enumerate}
\item \textbf{A validated simulator.} We reproduce a published analog-training result (Gokmen and Haensch 2020) with an
      independent device model, and we qualitatively check the offset-robustness trend of Rasch et al.\ 2024
      (Section~\ref{sec:validation}).
\item \textbf{A pre-registered comparison} of plain SGD and three device-tolerant algorithms on four published electrochemical
      device specifications and Widrow's memistor. There are three tasks (MNIST, a CIFAR-10 subset, and a character-level language
      model on human-written text) and five seeds, with confidence intervals (Section~\ref{sec:main}).
\item \textbf{A scale study} in which every arm gets its own learning rate at every width. It tests an earlier observation that
      the device gap grows with model width (Section~\ref{sec:scale}).
\item \textbf{A device specification sheet} built from one-at-a-time and pairwise sensitivity sweeps: the minimum levels, and
      the maximum asymmetry, nonlinearity, write noise, drift and reference error for a given accuracy gap at each model size
      (Section~\ref{sec:spec}).
\item \textbf{A system-level energy model} with periphery, built only from published figures and our simulated pulse counts. It
      compares against an H100 at its measured utilisation and separates what is measured from what is projected
      (Section~\ref{sec:energy}).
\item \textbf{A roadmap} of device experiments with pass criteria, array sizes, a hardware-in-the-loop protocol, and a bench
      replication of the 1960 memistor with safety notes (Section~\ref{sec:roadmap}, Appendix~\ref{app:bench}).
\end{enumerate}

Every experiment was pre-registered before it ran. The plan, every deviation, the simulator, its checks (each shown to fail on a
planted fault), every run record and the scripts that turn records into the numbers and figures of this paper are released
together (Section~\ref{sec:availability}). All results are simulations, and every device parameter carries a provenance label
(measured, fitted, claimed, estimated or assumed).
